% Procedencia: integral 2249, cap. 24, pp. 609--618; capitulo_24_dualidad_t.tex.
% Las fórmulas de dominio operatorio se explicitan sobre la realización ya construida.
\section{Dualidad T: coordenadas enteras, cambio de sección y equivalencia operatoria}
\label{iv:dualidad}

La descomposición aritmética de APP conserva el representante residual
y su cociente. La evolución TPK transporta además su orientación y
la memoria de los cruces de frontera. El intercambio de las dos
coordenadas exige conservar esta distinción y precisar el dominio
en que la operación vuelve a estar definida. Sobre ese dominio se
construye primero una involución aritmética; su forma cuadrática
admite después una representación por operadores autoadjuntos.

\subsection{Dominio estable del intercambio}

En la sección residual positiva,
\begin{equation}
 r_9^+(n)=1+((n-1)\bmod9),\qquad
 Q_9^+(n)=\frac{n-r_9^+(n)}9,
 \qquad n=r_9^+(n)+9Q_9^+(n).
 \label{iv:eq:division-positiva}
\end{equation}
Un intercambio puede situar un cociente arbitrario en la primera
coordenada. Por ello su dominio estable es
\(\mathscr D_0=\ZZ^2\times\RR_{>0}\), no solamente la sección
\(1\leq r\leq9\). La prolongación del dominio permite definir
una operación total; la admisibilidad de una pareja como lectura de
una ruta TPK conserva por separado sus reglas de selección.

\begin{definition}[Forma compacta normalizada e intercambio]
\label{iv:def:dualidad}
Para \(d=(m,w,R_0)\in\mathscr D_0\), sean
\begin{equation}
 E_{R_0}(m,w)=\frac{m^2}{R_0^2}+w^2R_0^2,
 \qquad
 \mathsf T_0(m,w,R_0)=(w,m,R_0^{-1}).
 \label{iv:eq:energia-t}
\end{equation}
Los símbolos \(m,w\) nombran las dos coordenadas enteras
transportadas; \(R_0\) es la escala positiva adimensional de
la forma. Su interpretación mediante momento, enrollamiento y
radio físico se establece por un mapa de realización posterior.
\end{definition}

\begin{theorem}[Involución y conservación de la forma cuadrática]
\label{iv:thm:dualidad-escalar}
La aplicación \(\mathsf T_0\) es una involución de
\(\mathscr D_0\), y
\[
 E_{R_0^{-1}}(w,m)=E_{R_0}(m,w).
\]
\end{theorem}
\begin{proof}
Dos intercambios restituyen \((m,w)\), y dos inversiones restituyen
\(R_0\). Además,
\[
 E_{R_0^{-1}}(w,m)=w^2R_0^2+\frac{m^2}{R_0^2}
 =E_{R_0}(m,w).
\]
\end{proof}

La equivalencia vale para cada radio y su recíproco. En \(R_0=1\),
el radio queda fijo y la operación actúa sobre una misma forma.
Una pareja concreta es fija por la involución completa precisamente
cuando, además, \(m=w\).

\subsection{Hamiltoniano compacto y dominio autoadjunto}

Sea \(\mathcal H_{\mathrm{lat}}=\ell^2(\ZZ^2)\), con base
ortonormal \(\{|m,w\rangle\}\). Definimos
\begin{align}
 H(R_0)|m,w\rangle&=E_{R_0}(m,w)|m,w\rangle,\label{iv:eq:hamiltoniano}\\
 D(H(R_0))&=\left\{\psi:\sum_{m,w\in\ZZ}
 E_{R_0}(m,w)^2|\psi_{m,w}|^2<\infty\right\}.\label{iv:eq:dominio-h}
\end{align}
Los vectores de soporte finito pertenecen al dominio y forman un
núcleo del operador diagonal. La forma cuadrática asociada tiene
dominio
\[
 D(H(R_0)^{1/2})=\left\{\psi:\sum_{m,w}
 E_{R_0}(m,w)|\psi_{m,w}|^2<\infty\right\}.
\]

\begin{proposition}[Autoadjunción y resolvente compacto]
\label{iv:prop:h-autoadjunto}
El operador \(H(R_0)\) es autoadjunto y no negativo. Sus
autovalores son \(E_{R_0}(m,w)\), contados con multiplicidad,
y \((H(R_0)+I)^{-1}\) es compacto.
\end{proposition}
\begin{proof}
Un operador de multiplicación por una sucesión real es simétrico
en \eqref{iv:eq:dominio-h}. Si \(\eta\) pertenece al dominio
del adjunto, probar la identidad del adjunto contra cada vector
\(|m,w\rangle\) obliga a que su imagen tenga coordenadas
\(E_{R_0}(m,w)\eta_{m,w}\). La condición de pertenecer a
\(\ell^2\) es exactamente \eqref{iv:eq:dominio-h}; por tanto
el adjunto tiene el mismo dominio y la misma acción. La no
negatividad se obtiene sumando los pesos no negativos.

Con \(c_{R_0}=\min(R_0^{-2},R_0^2)>0\), se tiene
\(E_{R_0}(m,w)\geq c_{R_0}(m^2+w^2)\). Fuera de cualquier
conjunto finito creciente, los coeficientes
\((1+E_{R_0}(m,w))^{-1}\) tienden uniformemente a cero.
Truncar el resolvente en dichos conjuntos produce operadores
de rango finito que convergen en norma, lo cual prueba su
compacidad. La base diagonal determina el espectro anunciado.
\end{proof}

\begin{theorem}[Conjugación unitaria, incluidos los dominios]
\label{iv:thm:t-unitaria}
\label{prop:dualidad-radios-k}
El intercambio
\(U_T|m,w\rangle=|w,m\rangle\) es unitario,
\(U_T^2=I\), y satisface
\begin{equation}
 U_TD(H(R_0))=D(H(R_0^{-1})),
 \qquad U_TH(R_0)U_T^{-1}=H(R_0^{-1}).
 \label{iv:eq:conjugacion-h}
\end{equation}
También transporta los dominios de las formas cuadráticas.
\end{theorem}
\begin{proof}
El intercambio es una permutación involutiva de la base
ortonormal. Para una sucesión \(\psi\), la identidad del
teorema \ref{iv:thm:dualidad-escalar} da
\[
 \sum_{m,w}E_{R_0^{-1}}(m,w)^2|(U_T\psi)_{m,w}|^2
 =\sum_{m,w}E_{R_0}(m,w)^2|\psi_{m,w}|^2.
\]
Así se obtiene la igualdad de dominios. Sobre cada coordenada,
\(U_TH(R_0)U_T^{-1}\) multiplica por
\(E_{R_0}(w,m)=E_{R_0^{-1}}(m,w)\). La misma suma, con
la primera potencia del peso, demuestra el resultado para las
formas cuadráticas.
\end{proof}

La igualdad escalar y la igualdad operatoria son dos realizaciones
de la misma involución. La segunda conserva los dominios que
permiten interpretar la primera como equivalencia espectral,
no sólo como una igualdad sobre vectores de soporte finito.

\subsection{Especialización al radio de la elipse de acción}

El radio \(R_\star\) determinado en la sección
\ref{iv:ellipse:section} fija la realización de esta colección que
corresponde a las representaciones angulares del estado. El símbolo
\(R_0\) permite expresar los resultados de transporte para toda la
colección positiva; la composición con la elipse evalúa esa colección en
\(R_0=R_\star\) y conserva también su imagen recíproca.

\begin{corollary}[Colección elíptica estable por dualidad]
\label{iv:cor:radio-especializado}
El conjunto
\[
 \mathscr D_{\mathrm{el}}
 =\ZZ^2\times\{R_\star,R_\star^{-1}\}
\]
es estable por \(\mathsf T_0\). Sus Hamiltonianos son autoadjuntos,
no negativos y de resolvente compacto. El intercambio \(U_T\)
transporta el dominio de \(H(R_\star)\) sobre el de
\(H(R_\star^{-1})\), incluidos los dominios de sus formas.
\end{corollary}
\begin{proof}
La positividad de \(R_\star\) y su determinación única se establecen
en la sección \ref{iv:ellipse:section}. La inversión permuta
\(R_\star\) y \(R_\star^{-1}\), mientras el intercambio preserva
\(\ZZ^2\). Se aplican, en estas dos fibras, la proposición
\ref{iv:prop:h-autoadjunto} y el teorema \ref{iv:thm:t-unitaria}.
Si \(R_\star=1\), las dos fibras coinciden; en cualquier otro caso
el operador entrelaza dos fibras distintas de la misma colección.
\end{proof}

La escala común de acción determina el área de la elipse; su cociente
de semiejes determina \(R_\star\). Esta separación permite conservar
el área y la escala común de acción al intercambiar los semiejes y obtener la reciprocidad
del radio. El intercambio de índices sobre \(\ell^2(\ZZ^2)\)
y los transportes del plano de acción actúan en espacios diferentes;
sus propiedades respectivas están demostradas en sus dominios propios.

\subsection{Cambio de representante y memoria del cociente}

La sección cero usa \(\bar r\in\{0,\ldots,8\}\). El cambio
que conserva el entero de \eqref{iv:eq:division-positiva} es
\begin{equation}
 C(r,Q)=(\bar r,Q+\mathbf1_{\{r=9\}}).
 \label{iv:eq:cambio-seccion}
\end{equation}
La modificación del cociente forma parte de la operación.
Para \(n=9\), las dos parejas son \((9,0)\) y \((0,1)\).
Sus energías no corregidas son \(81/R_0^2\) y \(R_0^2\).
Representar el mismo entero no implica, por tanto, igualdad de
esas dos inserciones en una forma que distingue las coordenadas.

La relación de equivalencia se conserva definiendo
\[
 L_9=\ZZ(9,-1),\qquad S(x_1,x_2)=(x_2,x_1),
 \qquad q_{R_0}(x)=\frac{x_1^2}{R_0^2}+x_2^2R_0^2.
\]
Para \(x\in\ZZ^2\), sea
\begin{equation}
 \mathcal E_{R_0}^{L_9}([x])
 =\min_{k\in\ZZ}q_{R_0}(x+k(9,-1)).
 \label{iv:eq:energia-cociente}
\end{equation}
Esta construcción opera sobre clases. La memoria que distingue
representantes individuales permanece en el estado TPK anterior
al cociente; la función de \eqref{iv:eq:energia-cociente}
determina la evaluación invariante de la clase elegida.

\begin{theorem}[Dualidad covariante de las secciones]
\label{iv:thm:dualidad-cociente}
La energía de \eqref{iv:eq:energia-cociente} está bien definida.
La transformación
\[
 \mathsf T_{\mathrm{cov}}([x]_{L_9},R_0)
 =([Sx]_{SL_9},R_0^{-1})
\]
es una biyección entre
\((\ZZ^2/L_9)\times\RR_{>0}\) y
\((\ZZ^2/SL_9)\times\RR_{>0}\), con inversa dada por el
segundo intercambio, y
\[
 \mathcal E_{R_0^{-1}}^{SL_9}([Sx])
 =\mathcal E_{R_0}^{L_9}([x]).
\]
\end{theorem}
\begin{proof}
Reemplazar \(x\) por \(x+k_0(9,-1)\) reindexa los candidatos
al mínimo. Éste existe: la función de \(k\) es cuadrática y
tiene coeficiente principal \(81/R_0^2+R_0^2>0\). El cambio
\((9,Q)\mapsto(0,Q+1)\) tiene diferencia en \(L_9\), por
lo que representa una misma clase. El intercambio transporta
la relación de equivalencia a \(SL_9\), y
\[
 \min_{\ell\in L_9}q_{R_0^{-1}}(Sx+S\ell)
 =\min_{\ell\in L_9}q_{R_0}(x+\ell).
\]
Finalmente, \(S^2=I\) recupera el cociente y el radio iniciales.
\end{proof}

La evaluación del mínimo puede efectuarse exactamente con dos
candidatos. Si \(x=(m,w)\), el mínimo real de la cuadrática
ocurre en
\[
 k_* =\frac{wR_0^4-9m}{81+R_0^4}.
\]
El mínimo entero se alcanza en \(\lfloor k_*\rfloor\) o en
\(\lceil k_*\rceil\). Esta fórmula no altera la equivalencia;
explicita su evaluación finita y permite comprobar el efecto
del cambio de sección sin descartar el cociente.

\subsection{Normalización de radio en la realización de cuerda}

Sea \(\alpha'>0\) una escala de la realización de cuerda.
Este símbolo designa la escala de cuerda y se mantiene distinto
de la constante de estructura fina \(\alpha\). Definimos
\begin{equation}
 \mathsf T_{\alpha'}(m,w,R)
 =\left(w,m,\frac{\alpha'}R\right),\qquad
 \mathcal N_{\alpha'}(m,w,R)
 =\left(m,w,\frac R{\sqrt{\alpha'}}\right).
 \label{iv:eq:normalizacion-cuerda}
\end{equation}

\begin{theorem}[Conjugación exacta por la escala]
\label{iv:thm:normalizacion-cuerda}
La normalización satisface
\[
 \mathcal N_{\alpha'}\mathsf T_{\alpha'}
 =\mathsf T_0\mathcal N_{\alpha'}.
\]
Si
\[
 E_{\mathrm{str}}(m,w,R)
 =\frac{m^2}{R^2}+\frac{w^2R^2}{\alpha'^2},
\]
entonces
\(E_{\mathrm{str}}=\alpha'^{-1}E\circ\mathcal N_{\alpha'}\).
Las combinaciones
\[
 p_L=\frac mR+\frac{wR}{\alpha'},\qquad
 p_R=\frac mR-\frac{wR}{\alpha'}
\]
se transforman mediante \(p_L\mapsto p_L\) y \(p_R\mapsto-p_R\).
\end{theorem}
\begin{proof}
Las dos composiciones de la primera igualdad dan
\((w,m,\sqrt{\alpha'}/R)\). Por otra parte,
\[
 E_{R/\sqrt{\alpha'}}(m,w)
 =\frac{\alpha'm^2}{R^2}+\frac{w^2R^2}{\alpha'}
 =\alpha'E_{\mathrm{str}}(m,w,R).
\]
Sustituir \((w,m,\alpha'/R)\) en las expresiones de
\(p_L,p_R\) da, respectivamente, las dos transformaciones
afirmadas.
\end{proof}

La conjugación fija qué significa expresar el mismo resultado en
una coordenatización dimensional. En la colección elíptica, los radios son
\[
 R=\sqrt{\alpha'}R_\star,\qquad
 R^\vee=\sqrt{\alpha'}R_\star^{-1}.
\]
Su producto es \(\alpha'\). La determinación interna del cociente
adimensional de semiejes y la escala de longitud al cuadrado de la
realización cumplen funciones distintas: la primera fija la forma;
la segunda expresa sus radios en una unidad dimensional.

La normalización anterior permite expresar el mismo resultado en
radio normalizado o físico. La asignación de una tensión y una
unidad a \(\alpha'\) pertenece al mapa de realización y tiene
su propio origen; no es un dato que intervenga en la producción
de los registros APP--TRIT--TPK.

\subsection{Relación con la transformación modular}

La inversión del radio tiene una representación precisa en el
eje imaginario del semiplano superior. Sean
\(\iota(R_0)=\mathrm iR_0^2\) y
\(\mathsf S(\tau)=-1/\tau\). Entonces
\begin{equation}
 \mathsf S\circ\iota(R_0)=\iota(R_0^{-1}).
 \label{iv:eq:semiconjugacion-modular}
\end{equation}
La igualdad resulta de
\(-1/(\mathrm iR_0^2)=\mathrm iR_0^{-2}\). Para la función
modular \(J=j-744\), cuya invariancia se conserva en el
desarrollo de Moonshine, se deduce
\(J(\mathrm iR_0^2)=J(\mathrm iR_0^{-2})\).

El mapa \(\iota\) expresa una semiconjugación de parámetros.
El intercambio \(U_T\) actúa en las coordenadas de
\(\ell^2(\ZZ^2)\); \(\mathsf S\) actúa en el semiplano
superior; el automorfismo dual de un orbifold actúa en sus
sectores graduados. La escritura de sus dominios permite
componer relaciones específicas entre ellos conservando la
información de cada representación. Los teoremas
\ref{iv:thm:t-unitaria}, \ref{iv:thm:dualidad-cociente}
y \ref{iv:thm:normalizacion-cuerda} establecen, respectivamente,
la conjugación unitaria, el cambio de sección residual y la
normalización de escala, con los dominios operatorios declarados.
